\newacronym{pic}{PIC}{Position Independent Code, позиционно независимый код}
\newacronym{got}{GOT}{Global Offset Table, таблица глобальных смещений}
\newacronym{plt}{PLT}{Procedure Linkage Table, таблица связывания процедур}
\newacronym{aslr}{ASLR}{Address Space Layout Randomization, рандомизация адресного пространства}
\newacronym{os}{ОС}{Операционная Система}
\newacronym{sw}{ПО}{Программное обеспечение}
\newacronym{vm}{ВМ}{Языковая Виртуальная Машина}
\newacronym{dfg}{ГПД}{Граф Потока Данных. Вершины графа моделируют исполнения инструкций (иногда называемые просто инструкциями), ребра - зависимости между инструкциями по значениям. С каждым ребром ассоциируются значение и локация, по которой данное значение располагается (также называемые значением и локацией ребра).}

\newglossaryentry{module}
{
    name={Модуль},
    description={Библиотека или приложение}
}

\newglossaryentry{unit-testing}{
    name={Юнит тестирование (unit testing)},
    description={Тестирование минимального фрагмента кода (функции или метода) в изоляции от других компонентов. Поведение внешних зависимостей имитируется с помощью объектов-подражателей (mock objects)}
}

\newglossaryentry{module-testing}{
    name={Модульное тестирование (module testing)},
    description={Тестирование группы связанных компонентов, образующих законченный модуль системы. Проверяет корректность взаимодействия внутри модуля и его интерфейсов}
}

\newglossaryentry{end-to-end-testing}{
    name={Сквозное тестирование (end to end testing)},
    description={Тестирование полного бизнес-сценария в условиях, максимально приближенных к реальным. Проверяет работу всех компонентов системы вместе, включая внешние зависимости}
}

\newglossaryentry{sink}
{
    name={Псевдо-инструкция},
    description={Вершина ГПД, не отвечающая реальной инструкции. Используется для отображения различных зависимостей на графе}
}

\newglossaryentry{addr_edge}
{
    name={Адресное ребро},
    description={Ребро ГПД, значение которого используется в качестве адреса}
}

\newglossaryentry{src_addr_edge}{
    name={Исходное адресное ребро},
    description={Ребро, на котором заканчивается построение обратного пути из адресного ребра на ГПД в алгоритме поиска адресов}
}

\newglossaryentry{trace}{
    name={Трасса исполнения},
    description={Набор упорядоченных записей о событиях, происходивших во время трассировочного запуска приложения}
}

\newglossaryentry{dfg_math}{
    name={$G = (V, E)$},
    description={ГПД, где: $\mathbf{V}$ - множество вершин графа, $\mathbf{E}$ - множество ребер графа},
    type=symbols
}

\newglossaryentry{edges_subsets}{
    name={$\mathbf{2^E}$},
    description={Множество всех подмножеств ребер ГПД},
    type=symbols
}

\makeglossaries
